\section{Discussion}
Mean fibroblast loading and the fraction of high-scoring fibroblasts were closely associated across the discovery specimens, and the association persisted after patient aggregation. This establishes that the mean conveys information about the prevalence of high scores among recovered fibroblasts. It does not make the measurements interchangeable. Low- and high-bin means varied across specimens, and some specimens contained no high-bin cells. An approximation that replaced these means with pooled values left a measurable residual. The defensible interpretation therefore includes both prevalence and within-bin score variation, conditional on the representation, gate and threshold used to define the measurement \citep{altman2006,buettner2021}.

The patient-unit extension connects the fibroblast interpretation directly to the available healing categories. Averaging programme scores within fibroblasts, counting high-scoring fibroblasts and counting fibroblasts among all cells all gave higher point estimates in healed patients, but each comparison retained considerable uncertainty. Reweighting repeated biopsies did not remove that uncertainty. A patient-level comparison of the entire loading distribution also retained uncertainty, while equal-cell sampling and every one-specimen-per-patient choice preserved the direction of the point contrast. These checks address selected measurement and sampling choices within the discovery material; they do not validate the readout in a new cohort. The appropriate implication is that the observed group direction warrants prospective evaluation of a fixed measurement, not that a threshold has been clinically validated or that a null biological effect has been demonstrated. Correcting the unit of analysis is necessary even when a sample-level contrast appears more favorable \citep{squair2021,zimmerman2021}.

The different uncertainty summaries answer different questions. Patient bootstrap intervals resample the observed outcome arms and inherit the limitations of their small empirical distributions \citep{efron1993}. Exhaustive label enumeration includes the observed allocation, whereas an add-one correction is commonly motivated for randomly sampled permutations \citep{phipson2010,hemerik2018}. Retaining the earlier convention alongside the exhaustive fraction makes that arithmetic transparent without making observational labels randomized. Likewise, Benjamini--Hochberg adjustment concerns a stated testing family under its dependence assumptions \citep{benjamini1995}; it cannot repair repeated specimens treated as independent patients. The retrospective power simulation substitutes the discovery effect and specimen spread into a hypothetical sampling model. It supplies no independent evidence about whether the uncertain patient association is real, absent or clinically important.

The cell atlas and expression profiles also clarify what is still missing experimentally. A loading enriched in stromal cells is not a direct count of an anatomically defined fibroblast population. Independent spatial or multiplex protein measurements in mapped biopsies would be needed to localize high-scoring cells and test whether their abundance changes within comparable tissue compartments \citep{foster2021,philippeos2018}. Perturbing a selected programme in an appropriate wound system could then ask whether it contributes to repair, rather than merely accompanying it \citep{mascharak2021}. These experiments address localization and function; neither can be supplied by adding more resamples to the present cell matrix.

This distinction has biological consequences. A high mean could accompany recruitment or recovery of a particular fibroblast population, increased expression within an existing population, or changes in several populations at once. Studies of dermal lineages and cross-tissue fibroblast organization show why these possibilities cannot be resolved by a small marker list alone \citep{driskell2013,muhl2020,buechler2021}. The enrichment of matrix and inflammatory genes supplies a useful description of the loading, but does not identify a unique lineage, a stable attractor or a causal repair mechanism. The broader separation between cell identity and cellular activity remains relevant when a continuous transcriptional programme is summarized at the tissue level \citep{kotliar2019}.

The threshold should be understood as an operational convention. It permits a reproducible high-state fraction and makes the denominator explicit, while assigning nearby scores on opposite sides of the cut to different bins. Information within those bins remains in the continuous score. This is one reason to retain both summaries rather than interpreting the binary fraction as a complete description of the programme \citep{altman2006}. The sensitivity to topic capacity also shows that a gene-associated coordinate need not have identical meaning across independently fitted representations. Neighborhood abundance methods could provide a complementary analysis of continuous state space, but they would introduce their own representation, neighborhood and replication choices \citep{dann2022}. They cannot retrospectively turn a descriptive threshold into proof of a discrete biological state.

The technical controls provide evidence about selected perturbations of this operational measurement. Most state labels survived count matching under fixed lineage assignments, and specimen fractions were stable after specified immune-transcript exclusions. Rerunning QC and the frozen gate on all released count columns reveals a different result: low depth removes many original cells, changes lineage labels among survivors and increases forced assignments without detected marker support. Stability conditional on a fixed or surviving population is therefore not stability of the whole tissue measurement. A relatively stable patient contrast can coexist with changed population membership and should not obscure that sensitivity. Neither stress test establishes cell-identity truth or excludes ambient RNA and doublets. Dedicated contamination and doublet procedures address different aspects of those processes \citep{young2020,wolock2019}. Dissociation responses can remain entangled with genuine injury responses \citep{vandenbrink2017}, and the similar attenuation of residual bimodality by immune and unrelated covariates limits attribution to one immune mechanism.

The difference between a tissue measurement and a healing predictor becomes clearer at the patient level. The healthy specimen with the highest fraction challenges disease exclusivity, while remaining compatible with a probabilistic association in a defined clinical population. Collapsing repeated specimens reduced the all-cell healing contrast and made its uncertainty explicit. This comparison concerns a different readout from the fibroblast-restricted fraction, so their effect estimates should not be treated as interchangeable measures of the same clinical association. More cells improve characterization of the sampled tissue but do not increase the number of independent outcomes \citep{squair2021,zimmerman2021}. The available healing categories also do not support a model of individual time to closure \citep{armstrong2017}.

The representation comparison remains an internal audit. Historical pooled-fold PCA and NMF coordinates are not directly comparable across fitted models, and their bootstrap contrasts cannot establish representation quality. The revised held-pair procedure uses the same fitted scoring rule for each healing/non-healing comparison and makes NMF component selection invariant to positive rescaling. It repairs that computational comparison without creating independent data: all pairs reuse the same 11 patients, and the frozen topic encoder and gene panel already used discovery expression. Overlapping training folds also share information \citep{bengio2004}. Exploring several thresholds, capacities and representations in the same cohort makes independent evaluation especially important, even when each individual comparison has a clearly specified computational rule \citep{cawley2010}. The uncertainty concerns the selection and transport of the readout as well as the arithmetic of a particular discrimination estimate.

Even reliable discrimination would leave other clinical questions unanswered. AUC describes ordering of outcomes, whereas calibration concerns whether predicted risks agree with observed risks in the population where the model is used \citep{steyerberg2010,vancalster2019}. The current loadings are expression summaries rather than calibrated clinical probabilities. Their value for a decision would also depend on the consequences of false reassurance and unnecessary intervention at clinically meaningful thresholds \citep{vickers2006}. This study does not fit such a decision model. Its contribution is to make the measurement and analysis unit sufficiently explicit that a later prognostic study can evaluate incremental value beyond routinely available clinical information.

The paired anatomical contrast demonstrates sensitivity to a known tissue difference, but it also emphasizes the importance of sampling context. Fibroblast organization varies across dermal compartments and with age \citep{driskell2013,soleboldo2020}. A detectable contrast between foot and forearm therefore supports responsiveness of the frozen representation without establishing specificity for wound healing. Spatial localization and an independently measured protein or transcript readout would help distinguish altered cellular abundance from altered activity within a compartment \citep{foster2021,kotliar2019}. Such measurements would address a different uncertainty from increasing the number of cells sequenced from the same biopsy.

External projection adds tissue contexts in which the frozen coordinate can be computed. It does not make the cohorts interchangeable. Different sampling designs, gene coverage, cell yields and compartment definitions can alter a summary before any difference in healing biology is considered. Benchmarking of single-cell integration demonstrates the importance of assessing technical alignment and conservation of biological variation together \citep{luecken2022}. Freezing the representation reduces one source of analytical flexibility, but does not ensure comparable measurement across source studies. The inconsistent immune associations and absence of independently adjudicated patient healing outcomes limit these cohorts to construct evaluation. The smallest projection, with two specimens, is chiefly a check of computational feasibility.

Future evaluation needs a defined outcome horizon, fixed sampling and readout, and a prespecified prognostic estimand. Prediction-model sample size also depends on outcome frequency and model complexity \citep{riley2019}. Equivalence requires a justified margin, not the smallest margin compatible with the observed data \citep{lakens2017}; a margin is not a prerequisite for every association or discrimination analysis. The present contribution is an executable measurement audit, not a new topic-model algorithm or biological law. It quantifies the limits of equating mean loading with prevalence, conditional stability with full-pipeline robustness, or specimen contrasts with patient outcomes. Computational reproducibility makes those limits inspectable; independent patients and orthogonal measurements are needed to establish clinical value.
